\section{Operator: Aritmetika, Perbandingan, Logika, dan Penugasan}

Operator adalah simbol yang melakukan operasi pada nilai (operand). JavaScript menyediakan beberapa kategori operator: aritmetika untuk perhitungan matematika, perbandingan untuk membandingkan nilai, logika untuk operasi boolean, dan penugasan untuk mengubah nilai variabel \cite{jsInfo}.

\begin{table}[htbp]
\centering
\begin{tabular}{|P{2.5cm}|P{4.5cm}|P{5cm}|}
\hline
\textbf{Kategori} & \textbf{Operator} & \textbf{Catatan} \\
\hline
Aritmetika & \code{+}, \code{-}, \code{*}, \code{/}, \code{\%}, \code{**} & \code{\%} modulo, \code{**} pangkat \\
\hline
Perbandingan & \code{==}, \code{===}, \code{!=}, \code{!==}, \code{<}, \code{>} & \code{===} tanpa konversi tipe \\
\hline
Logika & \code{\&\&} (AND), \code{||} (OR), \code{!} (NOT) & Short-circuit evaluation \\
\hline
Penugasan & \code{=}, \code{+=}, \code{-=}, \code{*=}, \code{/=} & Shorthand operasi + assign \\
\hline
Increment & \code{++}, \code{--} & Pre (\code{++x}) vs post (\code{x++}) \\
\hline
\end{tabular}
\caption{Kategori operator JavaScript}
\end{table}

\begin{webcode}[language=JavaScript, caption={Operator perbandingan dan logika}]
// == vs === (ketat)
console.log(5 == "5");    // true  (konversi tipe)
console.log(5 === "5");   // false (tipe berbeda)
console.log(null == undefined); // true (kasus khusus)

// Short-circuit evaluation
const name = userName || "Guest"; // jika userName falsy, pakai "Guest"
const display = isAdmin && "Admin Panel"; // jika true, kembalikan string

// Nullish coalescing (??)
const port = config.port ?? 3000; // hanya null/undefined yang diganti
\end{webcode}

\begin{konsep}
\textbf{Selalu Gunakan \code{===} (Strict Equality).} Operator \code{==} melakukan konversi tipe sebelum membandingkan, yang sering menghasilkan perilaku tak terduga (\code{0 == ""} adalah \code{true}). Operator \code{===} membandingkan tanpa konversi---lebih aman dan lebih prediktif. Satu-satunya pengecualian yang wajar: \code{value == null} untuk memeriksa null atau undefined sekaligus \cite{mdnJsGuide}.
\end{konsep}

